% Generated draft: requires mathematical and source review.
\documentclass[11pt]{article}
\usepackage[margin=25mm]{geometry}
\usepackage{fontspec}
\setmainfont{Latin Modern Roman}
\usepackage{amsmath,amssymb,mathtools,unicode-math}
\setmathfont{Latin Modern Math}
\usepackage{enumitem,needspace,xurl,hyperref}
\setlength{\parindent}{0pt}
\setlength{\parskip}{5pt}
\setlength{\emergencystretch}{3em}
\newcommand{\N}{\mathbb{N}}
\newcommand{\Z}{\mathbb{Z}}
\newcommand{\Q}{\mathbb{Q}}
\newcommand{\R}{\mathbb{R}}
\newcommand{\C}{\mathbb{C}}
\newcommand{\F}{\mathbb{F}}
\newcommand{\A}{\mathbb{A}}
\newcommand{\PP}{\mathbb{P}}
\newcommand{\E}{\mathbb{E}}
\newcommand{\eps}{\varepsilon}
\newcommand{\dd}{\,\mathrm{d}}
\begin{document}
\section*{UnsolvedMath ID: 30003409}
\textbf{Source problem code:} catalogue entry
\subsection{Definitions and mathematical statement}
Identify each cardinal \(\mu\) with its initial ordinal. A tree \((T,<_{T})\) is a set equipped with a strict partial order such that, for every \(t\in T\), the predecessor set \(\{s\in T:s<_{T}t\}\) is well ordered by \(<_{T}\). Define
\[
\operatorname{ht}_{T}(t)=\operatorname{otp}(\{s\in T:s<_{T}t\},<_{T}),\qquad
T_{\alpha}=\{t\in T:\operatorname{ht}_{T}(t)=\alpha\},
\]
and
\[
\operatorname{ht}(T)=\sup\{\operatorname{ht}_{T}(t)+1:t\in T\}.
\]
A branch is a downward closed subset of \(T\) linearly ordered by \(<_{T}\). For a tree of height \(\mu\), a branch \(B\) is cofinal if
\[
(\forall\alpha<\mu)(\exists t\in B)\;\alpha\leq\operatorname{ht}_{T}(t).
\]
Using the explicit meaning of Aronszajn tree in the supplied record, a cardinal \(\mu\) carries such a tree if there exists a tree \((T,<_{T})\) satisfying
\[
\operatorname{ht}(T)=\mu,\qquad
(\forall\alpha<\mu)\;|T_{\alpha}|<\mu,
\qquad
T\text{ has no cofinal branch}.
\]
\textbf{Problem.} Characterize the cardinals \(\mu\) carrying a tree with these properties.

The supplied formulation imposes no explicit infinitude, uncountability, or regularity hypothesis on \(\mu\); none is added here. Whether a narrower conventional domain was intended remains unspecified. The requested form of characterization, including whether it concerns existence in a fixed model of ZFC or relative consistency across models, is also unspecified.
\subsection{Short English statement}
Determine which cardinals \(\mu\) admit a tree of height \(\mu\) whose individual levels all have size less than \(\mu\), but which has no branch reaching cofinally through its levels. The supplied record does not restrict \(\mu\) to regular uncountable cardinals or specify the intended kind of classification.
\subsection{Sources}
\begin{itemize}
\item Supplied local UnsolvedMath registry record, permanent ID \(30003409\), title \emph{Cardinals Carrying Aronszajn Trees}, source problem code \texttt{catalogue entry}, in the Oberwolfach Reports Open Problems collection. The record gives proposed year 2017 and supplies no named proposer.
\item Supplied source citation: \emph{Set Theory (2018)}, Oberwolfach Reports, pp.~527--589, DOI \texttt{10.4171/owr/2017/11}. Supplied source link: \url{https://doi.org/10.4171/owr/2017/11}. The background describes Oberwolfach Reports as documenting mathematical workshops and their open-problem discussions.
\item Supplied literature source, labelled ``OWR Set Theory report'': \url{https://ems.press/content/serial-article-files/46674}.
\item Supplied literature source, labelled ``Diagonal supercompact Radin forcing'': \url{https://doi.org/10.1016/j.apal.2020.102828}.
\item Literature assessment dated exactly \texttt{2026-08-21}: the supplied curation labels the problem open and states that there is no ZFC classification of the cardinals carrying Aronszajn trees, that tree property results and forcing constructions provide substantial partial information, and that the general question remains independent and unresolved. This is supplied dated curation derived from the OWR record, not an independently verified assessment. Its time-sensitive status and cited claims require independent verification.
\item The registry labels the statement \texttt{corrected\_verified} and says that mathematical content was preserved while notation, typography, and sentence structure were repaired for publication. These are supplied registry assertions, not independent mathematical verification of this draft.
\end{itemize}
\end{document}
